% ===========================================================================
\chapter{What we are going to build}\label{ch:intro}

A family is a handful of facts. Alice is Bob's parent; Bob is Carol's
parent; Carol is Emma's parent. A handful of facts, written down, imply a
great many more: Alice is Emma's ancestor, Carol and Dan are cousins, Bob
and Beth are siblings. Nobody wrote those down, but anyone who knows what
``ancestor'' means can work them out.

\emph{Datalog} is a small language for exactly this. A program states some
facts and some rules, and the rules say how new facts follow from old ones:

\begin{shell}
parent(alice, bob).
parent(bob, carol).

ancestor(X, Y) :- parent(X, Y).
ancestor(X, Z) :- parent(X, Y), ancestor(Y, Z).
\end{shell}

The first rule says: if \code{X} is a parent of \code{Y}, then \code{X} is
an ancestor of \code{Y}. The second says: if \code{X} is a parent of some
\code{Y}, and \code{Y} is an ancestor of \code{Z}, then \code{X} is an
ancestor of \code{Z}. Then we ask questions:

\begin{shell}
?- ancestor(alice, Who).
\end{shell}

and the answer is every \code{Who} that makes it true.

By the end of this tutorial you will have written, in Nex, a Datalog engine
that reads programs like this one, works out every fact they imply, and
answers questions about them. It also explains its answers: ask it
\emph{why} Carol and Dan are cousins, and it shows the proof, step by step,
down to the facts you gave it.

\laterout{datalog.nex family.dl "why cousin(carol, dan)."}{why_cousin}

\section{Why Datalog}

Datalog looks like a toy, and it is small enough to build in a few hundred
lines. It is also used in earnest. Query optimisers and database engines
evaluate recursive queries the way we will. Program analysers written in
Datalog find bugs in millions of lines of code, by stating facts about a
program (``this variable may point to that object'') and rules about how
those facts combine. Access-control systems decide who may read what with
rules much like the ones in Chapter~\ref{ch:negation}. In every case the
appeal is the same: you say \emph{what} follows from what, and the engine
works out \emph{how}.

That last part is where the interesting programming is. Working out
everything a program implies is a search that has to stop, has to find
everything, and should not do the same work twice. Those three
requirements are also three contracts, and this tutorial writes them down
as contracts and lets Nex check them while the engine runs.

\section{The plan}

\begin{description}[style=nextline, leftmargin=1.2em, itemsep=0.4em]
  \item[Reading Datalog.] The data model of terms, atoms and rules
    (Chapter~\ref{ch:terms}), a lexer that cuts the text into tokens
    (Chapter~\ref{ch:lexer}), and a parser that builds the model from them
    (Chapter~\ref{ch:parser}).
  \item[Facts and how rules use them.] Matching a pattern like
    \code{parent(X, bob)} against facts (Chapter~\ref{ch:matching}), a
    database to keep the facts in (Chapter~\ref{ch:database}), and joining
    the conditions of a rule (Chapter~\ref{ch:joins}).
  \item[Working out everything.] Applying the rules until nothing new
    appears, which is called reaching a \emph{fixpoint}
    (Chapter~\ref{ch:fixpoint}); the rules that make this possible
    (Chapter~\ref{ch:safety}); answering questions
    (Chapter~\ref{ch:queries}); and doing it much faster
    (Chapter~\ref{ch:seminaive}).
  \item[More expressive programs.] Negation, and why it needs the program
    to be evaluated in layers (Chapter~\ref{ch:negation}).
  \item[Speed and explanations.] An index (Chapter~\ref{ch:indexing}),
    and remembering how each fact was derived so that the engine can say
    why it holds (Chapter~\ref{ch:explanations}).
  \item[The finished program.] A command line (Chapter~\ref{ch:commandline}),
    and where to go next (Chapter~\ref{ch:next}).
\end{description}

Every line of the engine appears in this tutorial, in the order you would
write it, and every output shown comes from running the code exactly as
printed.
